\documentclass[11pt]{article}
\usepackage[utf8]{inputenc}
\usepackage[T1]{fontenc}
\usepackage{amsmath,amssymb,mathtools}
\usepackage{geometry}
\usepackage{booktabs}
\usepackage{listings}
\usepackage{xcolor}
\geometry{margin=1in}

\lstdefinestyle{python}{
  language=Python,
  basicstyle=\ttfamily\small,
  keywordstyle=\color{blue!70!black},
  stringstyle=\color{green!40!black},
  commentstyle=\color{gray!70!black},
  showstringspaces=false,
  breaklines=true,
  frame=single
}

\title{Implementation Note: Text Cross-Attention for the MDM Shortcut Predictor}
\author{}
\date{}

\begin{document}
\maketitle

\section{Goal}

For HumanML3D, the shortcut predictor should not be treated like an
image-latent predictor. Its source activation is a motion sequence, and the
conditioning signal is text. The proposed default predictor block is therefore
\[
  h = \mathrm{Conv1D}_{\mathrm{time}}(x),\qquad
  h = h + \mathrm{SelfAttn}(h),\qquad
  h = h + \mathrm{CrossAttn}(h, C),\qquad
  h = h + \mathrm{MLP}(h),
\]
where motion tokens query a context set $C$ containing the MDM condition token
$z_{tk}$ and, when available, CLIP text tokens.

This change is for the auxiliary depth-shortcut predictor. Evaluation metrics
such as FID, matching score, R-precision, and diversity must still sample from
the full MDM backbone, not from the shortcut predictor.

\section{Tensor convention}

The predictor operates on hidden activations after MDM has embedded the motion:
\[
  S_a \in \mathbb{R}^{B\times T\times D},
\]
where $B$ is batch size, $T$ is number of motion frames/tokens, and $D$ is the
MDM hidden dimension, e.g. $D=512$ for the HumanML3D setting.

The context tokens are
\[
  C = [z_{tk}; c_1;\ldots;c_M]\in \mathbb{R}^{B\times (1+M)\times D}.
\]
Here $z_{tk}$ is the original MDM conditioning token created from timestep and
text:
\[
  z_{tk}=f_t(t)+f_c(c),
\]
with the same classifier-free random text mask used by the backbone during
training. The tokens $c_1,\ldots,c_M$ are optional projected CLIP text tokens.
If the current MDM implementation only exposes the pooled CLIP embedding, use
\[
  C = [z_{tk}; f_{\mathrm{pool}}(c)].
\]
If only $z_{tk}$ is available, start with $C=[z_{tk}]$; this is still better
than giving the predictor no text/timestep context.

\section{Cross-attention module}

Use pre-norm residual cross-attention. Initialize the residual gate at zero so
the new path starts as an identity mapping and does not immediately disturb the
existing predictor.

\begin{lstlisting}[style=python]
import torch
import torch.nn as nn


class MotionTextCrossAttention(nn.Module):
    def __init__(
        self,
        dim: int,
        context_dim: int | None = None,
        num_heads: int = 8,
        dropout: float = 0.0,
        gated: bool = True,
    ):
        super().__init__()
        context_dim = context_dim or dim
        self.motion_norm = nn.LayerNorm(dim)
        self.context_norm = nn.LayerNorm(context_dim)
        self.context_proj = (
            nn.Identity() if context_dim == dim else nn.Linear(context_dim, dim)
        )
        self.attn = nn.MultiheadAttention(
            embed_dim=dim,
            num_heads=num_heads,
            dropout=dropout,
            batch_first=True,
        )
        self.gate = nn.Parameter(torch.zeros(())) if gated else None

    def forward(
        self,
        motion_tokens: torch.Tensor,          # [B, T, D]
        context_tokens: torch.Tensor,         # [B, S, D_ctx]
        context_padding_mask: torch.Tensor | None = None,  # [B, S], True = mask
    ) -> torch.Tensor:
        q = self.motion_norm(motion_tokens)
        ctx = self.context_proj(self.context_norm(context_tokens))
        out, _ = self.attn(
            query=q,
            key=ctx,
            value=ctx,
            key_padding_mask=context_padding_mask,
            need_weights=False,
        )
        if self.gate is not None:
            out = torch.tanh(self.gate) * out
        return motion_tokens + out
\end{lstlisting}

\section{Predictor block}

The recommended HumanML3D predictor block is temporal Conv1D followed by
self-attention, cross-attention, and MLP. Conv1D runs along the frame axis.

\begin{lstlisting}[style=python]
class TemporalConvSelfCrossBlock(nn.Module):
    def __init__(self, dim: int, num_heads: int = 8, mlp_ratio: float = 4.0):
        super().__init__()
        self.conv_norm = nn.LayerNorm(dim)
        self.temporal_conv = nn.Sequential(
            nn.Conv1d(dim, dim, kernel_size=3, padding=1, groups=1),
            nn.GELU(),
            nn.Conv1d(dim, dim, kernel_size=3, padding=1, groups=1),
        )

        self.self_norm = nn.LayerNorm(dim)
        self.self_attn = nn.MultiheadAttention(
            dim, num_heads=num_heads, batch_first=True
        )

        self.cross_attn = MotionTextCrossAttention(
            dim=dim,
            context_dim=dim,
            num_heads=num_heads,
            gated=True,
        )

        hidden = int(dim * mlp_ratio)
        self.mlp = nn.Sequential(
            nn.LayerNorm(dim),
            nn.Linear(dim, hidden),
            nn.GELU(),
            nn.Linear(hidden, dim),
        )

    def forward(self, h, context_tokens, context_padding_mask=None):
        # Temporal local smoothing: [B, T, D] -> [B, D, T] -> [B, T, D]
        u = self.conv_norm(h).transpose(1, 2)
        u = self.temporal_conv(u).transpose(1, 2)
        h = h + u

        u, _ = self.self_attn(
            self.self_norm(h),
            self.self_norm(h),
            self.self_norm(h),
            need_weights=False,
        )
        h = h + u

        h = self.cross_attn(h, context_tokens, context_padding_mask)
        h = h + self.mlp(h)
        return h
\end{lstlisting}

\section{Context builder}

The predictor should receive the same conditioning state as the backbone. Do
not use unmasked text in the predictor while the backbone receives masked text;
that creates a classifier-free guidance mismatch.

\begin{lstlisting}[style=python]
def build_predictor_context(
    timestep_cond: torch.Tensor,          # z_tk or time+text cond, [B, D]
    pooled_text: torch.Tensor | None,     # [B, D_text]
    text_tokens: torch.Tensor | None,     # [B, M, D_text]
    text_token_mask: torch.Tensor | None, # [B, M], True = mask
    pooled_text_proj: nn.Module,
    text_token_proj: nn.Module,
):
    tokens = [timestep_cond[:, None, :]]
    masks = [
        torch.zeros(
            timestep_cond.shape[0],
            1,
            dtype=torch.bool,
            device=timestep_cond.device,
        )
    ]

    if text_tokens is not None:
        tokens.append(text_token_proj(text_tokens))
        masks.append(text_token_mask)
    elif pooled_text is not None:
        tokens.append(pooled_text_proj(pooled_text)[:, None, :])
        masks.append(
            torch.zeros(
                timestep_cond.shape[0],
                1,
                dtype=torch.bool,
                device=timestep_cond.device,
            )
        )

    context_tokens = torch.cat(tokens, dim=1)
    context_padding_mask = torch.cat(masks, dim=1)
    return context_tokens, context_padding_mask
\end{lstlisting}

\section{Full predictor forward path}

For a sampled depth pair $(a,b)$, the predictor keeps the existing
direction--magnitude objective. Only the internal conditioning block changes.

\begin{lstlisting}[style=python]
def forward(
    self,
    source_act,          # S_a, [B, T, D]
    source_layer,        # a
    target_layer,        # b
    timestep_cond,       # z_tk or MDM emb, [B, D]
    pooled_text=None,
    text_tokens=None,
    text_token_mask=None,
):
    context, context_mask = build_predictor_context(
        timestep_cond=timestep_cond,
        pooled_text=pooled_text,
        text_tokens=text_tokens,
        text_token_mask=text_token_mask,
        pooled_text_proj=self.pooled_text_proj,
        text_token_proj=self.text_token_proj,
    )

    h = self.input_proj(source_act)
    h = h + self.layer_pair_embed(source_layer, target_layer)[:, None, :]
    h = h + self.timestep_proj(timestep_cond)[:, None, :]

    for block in self.blocks:
        h = block(h, context, context_mask)

    direction = self.direction_head(h)       # [B, T, D]
    delta_mag = self.magnitude_head(h)       # [B, T, 1] or pooled [B, 1]
    return direction, delta_mag
\end{lstlisting}

\section{Recommended flags}

For HumanML3D, make the cross-attention path the default for the
\texttt{hybrid\_deep\_12pct} predictor:

\begin{center}
\begin{tabular}{ll}
\toprule
Switch & Default \\
\midrule
\texttt{--shortcut-predictor} & \texttt{hybrid\_deep\_12pct} \\
\texttt{--shortcut-predictor-use-timestep} & on \\
\texttt{--shortcut-predictor-use-text-cross-attn} & on \\
\texttt{--shortcut-predictor-text-context} & \texttt{ztk\_plus\_pooled} initially \\
\texttt{--shortcut-predictor-cross-attn-heads} & $8$ \\
\texttt{--shortcut-predictor-cross-attn-gated} & on \\
\bottomrule
\end{tabular}
\end{center}

If CLIP token embeddings are exposed later, switch
\texttt{--shortcut-predictor-text-context} from \texttt{ztk\_plus\_pooled} to
\texttt{ztk\_plus\_tokens}.

\section{Implementation locations}

The minimal patch should touch these components:

\begin{itemize}
\item predictor module: add \texttt{MotionTextCrossAttention} and replace the
  pure CNN/MLP-only block with \texttt{TemporalConvSelfCrossBlock};
\item MDM forward path: return or expose the already computed $z_{tk}$ condition
  used by the backbone;
\item text encoder path: optionally expose CLIP token embeddings. This is not
  required for the first version if pooled text is already available;
\item training loop: pass $z_{tk}$ and optional text context into the shortcut
  predictor for direct, bootstrap, output-distill, and private losses;
\item args/config: add the cross-attention flags above and enable them by
  default for HumanML3D.
\end{itemize}

\section{Sanity checks}

Before running a full 200k-step job:

\begin{itemize}
\item Check shape: source activation and predicted direction must both be
  $[B,T,D]$.
\item Check classifier-free consistency: when text is masked for the backbone,
  the predictor context must use the masked text condition too.
\item Check cold-start behavior: with cross-attention gate initialized at zero,
  the first-step loss should match the previous predictor within noise.
\item Check non-finite counters: \texttt{output\_distill\_nonfinite} should stay
  at zero.
\item Check evaluation path: FID and R-precision must call the normal MDM
  sampling function, not the shortcut predictor.
\item Check parameter budget: keep the predictor near $10$--$12\%$ of the MDM
  backbone by controlling number of blocks, hidden width, heads, and MLP ratio.
\end{itemize}

\section{Expected effect}

Temporal Conv1D provides the local smooth-motion bias, self-attention captures
long-range motion structure, and cross-attention lets each motion token query
the actual text/timestep condition. This is a better HumanML3D inductive bias
than a predictor that only sees source activations and scalar timestep/layer
ids, because the shortcut target depends on the action semantics encoded in
text.

\end{document}
